\documentclass[../../main.tex]{subfiles}
\graphicspath{{../../image/}{../../figure/}}
\begin{document}

\section{Euler 方程组}\label{section:ch04-Euler方程}

支配一维、二维或三维等熵气体流动的 Euler 方程组\cref{chapter:ch03-双曲平衡律方程组}，提供了双曲守恒律组最原初、也仍然最重要的范例。它们长期作为整个类的范式，并拥有庞大的文献，既讨论与其特殊结构相联系的性质，也讨论其他双曲守恒律组所共有的一般性质。为给本书其余部分要关注的问题搭好舞台，本节中我们展示三维 Euler 方程组的一组代表性性质，它们来自新近或较早的研究工作。分析在此仅作勾勒--细节可查阅第 4.9 小节所引文献。


取零体积力 $b=0$，把\cref{chapter:ch03-双曲平衡律方程组}（3.3.36）写成正则形式的方程组，以质量密度 $\rho$ 与动量密度 $m=\rho v$ 为状态变量：

\begin{align}
\begin{cases}
\partial_{t}\rho+\operatorname{div}m^{\top}=0\\
\partial_{t}m+\operatorname{div}(\rho^{-1}mm^{\top})+\operatorname{grad}p(\rho)=0.
\end{cases}
\label{align:ch04-Euler}
\end{align}

回顾\cref{chapter:ch03-双曲平衡律方程组}，方程组\eqref{align:ch04-Euler}配有熵--熵流对

\begin{align}
\eta=\rho\varepsilon(\rho)+\tfrac{1}{2}\rho^{-1}|m|^{2},\qquad Q=\eta\rho^{-1}m+p(\rho)\rho^{-1}m,
\label{align:ch04-Euler熵}
\end{align}
其中 $\varepsilon'(\rho)=\rho^{-2}p(\rho)$。

在古典解的范围内，可合并\eqref{align:ch04-Euler}中两个方程，把方程组写成以 $(\rho,v)$ 为状态变量的形式
\begin{align}
\begin{cases}
\partial_{t}\rho+(v\cdot\operatorname{grad})\rho+\rho\operatorname{div}v=0\\
\partial_{t}v+(v\cdot\operatorname{grad})v+p'(\rho)\operatorname{grad}\rho=0.
\end{cases}
\label{align:ch04-Euler非守恒}
\end{align}

方程组\eqref{align:ch04-Euler}、\eqref{align:ch04-Euler非守恒}是双曲的，且只要 $p'(\rho)>0$，熵 $\eta$ 是 $(\rho,m)$ 的一致凸函数。本节通篇为简单起见假设气体是理想气体，状态方程为\cref{chapter:ch03-双曲平衡律方程组}（2.5.31），即 $p=\kappa\rho^{\gamma}$，其中按基本分子动理论 $\gamma\in(1,5/3]$。声速为 $c(\rho)=[\kappa\gamma\rho^{\gamma-1}]^{1/2}$。因而我们的方程组当 $\rho>0$ 时双曲，但双曲性在 $\rho=0$ 处失效。真空的出现使 Euler 方程组的分析复杂化。


可通过把 $\rho$ 替换为新状态变量 $\omega=\beta^{-1}c(\rho)$，其中 $\beta=\tfrac{\gamma-1}{2}$，从而在包括真空在内的整个物理范围上恢复双曲性，于是\eqref{align:ch04-Euler非守恒}变为
\begin{align}
\begin{cases}
\partial_{t}\omega+(v\cdot\operatorname{grad})\omega+\beta\omega\operatorname{div}v=0\\
\partial_{t}v+(v\cdot\operatorname{grad})v+\beta\omega\operatorname{grad}\omega=0,
\end{cases}
\label{align:ch04-Euler对称双曲}
\end{align}
即使当 $\omega=0$ 时这也是对称双曲组。正如我们将在\cref{chapter:ch05-熵与经典解稳定性}中看到，这保证了：当初值 $(\omega_{0},v_{0})$ 属于足够高阶 Sobolev 空间且 $\omega_{0}\geqslant 0$ 时，\eqref{align:ch04-Euler对称双曲}的 Cauchy 问题在某极大时间区间 $[0,T)$ 上存在唯一古典解 $(\omega,v)$。进而，若 $T$ 有限，则 $\omega$ 和/或 $v$ 的某导数必在 $t\to T$ 时爆破。显然，解 $(\omega,v)$ 在同一时间区间 $[0,T)$ 上诱导出\eqref{align:ch04-Euler}与\eqref{align:ch04-Euler非守恒}的古典解 $(\rho,m)$ 与 $(\rho,v)$。


我们现采用与第 4.2 小节中 Burgers 方程精神相似的论证，来证明上述古典解的寿命 $T$ 一般是有限的，这是质量囚禁的结果。考虑\eqref{align:ch04-Euler}的 Cauchy 问题，初值 $\rho(x,0)=\rho_{0}=\text{常数}$（对 $x\in\mathbb{R}^{3}$），$m(x,0)$ 支于单位球内，当 $|x|\geqslant 1$ 时 $m(x,0)=0$。设在某时间区间 $[0,T)$ 上存在古典解 $(\rho(x,t),m(x,t))$。方向 $\nu\in\mathbb{S}^{2}$ 上的快特征速度为 $v\cdot\nu+c(\rho)$。于是由\cref{theorem:ch04-影响区域}，对任一 $t\in[0,T)$ 与 $|x|\geqslant r(t)$，$\rho(x,t)=\rho_{0}$，$m(x,t)=0$，其中 $r(t)=1+c(\rho_{0})t$。特别地，由\eqref{align:ch04-Euler}的第一式，
\begin{align}
\int_{|x|<r(t)}[\rho(x,t)-\rho_{0}]\,\mathrm{d}x=0,\qquad 0\leqslant t<T.
\label{align:ch04-质量守恒}
\end{align}

我们监测加权径向动量的演化
\begin{align}
M(t)=\int_{|x|<r(t)}x\cdot m(x,t)\,\mathrm{d}x.
\label{align:ch04-径向动量}
\end{align}
对 $t$ 微分\eqref{align:ch04-径向动量}，通过\eqref{align:ch04-Euler}第二式把时间导数 $\partial_{t}m$ 表为空间导数，分部积分，得
\begin{align}
\dot{M}(t)=\int_{|x|<r(t)}\left[\rho^{-1}|m|^{2}+3\kappa(\rho^{\gamma}-\rho_{0}^{\gamma})\right]\,\mathrm{d}x.
\label{align:ch04-Mdot}
\end{align}
由于 $\gamma>1$，\eqref{align:ch04-质量守恒}与 Jensen 不等式蕴含
\begin{align}
\int_{|x|<r(t)}(\rho^{\gamma}-\rho_{0}^{\gamma})\,\mathrm{d}x\geqslant 0.
\label{align:ch04-Jensen}
\end{align}
进而，由\eqref{align:ch04-径向动量}、\eqref{align:ch04-质量守恒}与 Schwarz 不等式，
\begin{align}
M^{2}(t)\leqslant\int_{|x|<r(t)}\rho|x|^{2}\,\mathrm{d}x\int_{|x|<r(t)}\rho^{-1}|m|^{2}\,\mathrm{d}x\leqslant\frac{4\pi}{3}\rho_{0}r^{5}(t)\int_{|x|<r(t)}\rho^{-1}|m|^{2}\,\mathrm{d}x.
\label{align:ch04-Schwarz}
\end{align}
把\eqref{align:ch04-Mdot}与\eqref{align:ch04-Jensen}、\eqref{align:ch04-Schwarz}结合，得到微分不等式
\begin{align}
\dot{M}(t)\geqslant\frac{3}{4\pi\rho_{0}}[1+c(\rho_{0})t]^{-5}M^{2}(t).
\label{align:ch04-M微分}
\end{align}
经初等积分，并回顾 $c(\rho_{0})=[\kappa\gamma\rho_{0}^{\gamma-1}]^{1/2}$，我们得到：若
\begin{align}
M(0)>\frac{16\pi}{3}(\kappa\gamma)^{1/2}\rho_{0}^{\frac{\gamma+1}{2}},
\label{align:ch04-M临界}
\end{align}
则 $M(t)$ 将在有限时间内爆破。因而，大初值等熵气体动力学方程组 Cauchy 问题的古典解一般在有限时间内崩溃。

文献（第 4.9 小节所引）中报道的更精细分析表明，上述论证可推广以确立非等熵气体动力学中古典解的崩溃，甚至当初值未必很大时亦然。

上述论证验证了 Euler 方程组\eqref{align:ch04-Euler} Cauchy 问题古典解的寿命一般有限，但并未精确指出灾难在何时、以何种方式发生。典型地，古典解作为波破碎的结果而爆破；但在转向这一问题之前，我们先讨论在真空存在时质量囚禁的另一表现。

考虑方程组\eqref{align:ch04-Euler}以初值 $(\rho_{0},m_{0})$、$m_{0}=\rho_{0}v_{0}$ 的 Cauchy 问题，初值在 $\mathbb{R}^{3}$ 的有界子集 $\Omega$ 外为零。设在某时间区间 $[0,T)$ 上存在（一般为弱）解 $(\rho,m)$，它相对于熵--熵流对\eqref{align:ch04-Euler熵}满足熵可容性条件\eqref{align:ch04-熵积分tau}。气体将以有限速度弥散入真空，因而对每一 $t\in[0,T)$，$(\rho,m)$ 支于某有界集 $\Omega_{t}$。我们将借助下面这一有趣估计来估计 $\Omega_{t}$ 的大小。

从恒等式出发
\begin{align}
\tfrac{1}{2}\rho|tv-x|^{2}+t^{2}\rho\varepsilon=\tfrac{1}{2}|x|^{2}\rho-tx\cdot m+t^{2}\eta(\rho,m).
\label{align:ch04-恒等式}
\end{align}
经冗长但直接的计算，利用\eqref{align:ch04-Euler}与熵不等式，推出
\begin{align}
\partial_{t}\left(\tfrac{1}{2}\rho|tv-x|^{2}+t^{2}\rho\varepsilon\right)+\operatorname{div}\left((\tfrac{1}{2}\rho|tv-x|^{2}+t^{2}\rho\varepsilon)v+tp[tv-x]\right)\leqslant t[2\rho\varepsilon-3p].
\label{align:ch04-弥散不等式}
\end{align}
特别地，当 $(\rho,m)$ 为古典解时，\eqref{align:ch04-弥散不等式}取等式。

对理想气体\cref{chapter:ch03-双曲平衡律方程组}（2.5.31），$2\rho\varepsilon-3p=(5-3\gamma)\rho\varepsilon$。于是在 $\mathbb{R}^{3}$ 上积分\eqref{align:ch04-弥散不等式}得微分不等式

\begin{align}
\frac{\mathrm{d}}{\mathrm{d}t}\int_{\mathbb{R}^{3}}\left(\tfrac{1}{2}\rho|tv-x|^{2}+t^{2}\rho\varepsilon\right)\,\mathrm{d}x\leqslant\frac{5-3\gamma}{t}\int_{\mathbb{R}^{3}}\left(\tfrac{1}{2}\rho|tv-x|^{2}+t^{2}\rho\varepsilon\right)\,\mathrm{d}x,
\label{align:ch04-弥散微分}
\end{align}
由此
\begin{align}
\int_{\mathbb{R}^{3}}\rho\varepsilon\,\mathrm{d}x\leqslant Ct^{3(1-\gamma)}.
\label{align:ch04-内能衰减}
\end{align}
\eqref{align:ch04-内能衰减}的物理解释是：气体膨胀把内能转化为动能，结果内能衰减。

由 Hölder 不等式，
\begin{align}
\int_{\mathbb{R}^{3}}\rho\,\mathrm{d}x\leqslant\left(\int_{\mathbb{R}^{3}}\rho^{\gamma}\,\mathrm{d}x\right)^{1/\gamma}|\Omega_{t}|^{\frac{1-\gamma}{\gamma}}\leqslant Ct^{\frac{3(1-\gamma)}{\gamma}}|\Omega_{t}|^{\frac{1-\gamma}{\gamma}},
\label{align:ch04-Holder}
\end{align}
其中 $|\Omega_{t}|$ 表示 $\Omega_{t}$ 的体积。\eqref{align:ch04-Holder}左端是正常数，即气体的守恒质量。因而 $|\Omega_{t}|\geqslant At^{3}$，其中 $A$ 为某正常数。我们现讨论古典解能否在长期内以上述速率维持 $|\Omega_{t}|$ 的增长。

于是设上述解光滑，此时 $\Omega$ 与 $\Omega_{t}$ 为开集，其光滑边界记为 $\Gamma$ 与 $\Gamma_{t}$。在 $\Gamma_{t}$ 上，\eqref{align:ch04-Euler非守恒}与\eqref{align:ch04-Euler对称双曲}化为输运方程
\begin{align}
\partial_{t}v+(v\cdot\operatorname{grad})v=0.
\label{align:ch04-界面输运}
\end{align}
因而，对任一 $x\in\Gamma$，$v$ 沿直线 $y=x+tv_{0}(x)$ 保持常数 $v_{0}(x)$，且
\begin{align}
\Gamma_{t}=\{y\in\mathbb{R}^{3}:y=x+tv_{0}(x),\ x\in\Gamma\},
\label{align:ch04-界面Gamma_t}
\end{align}
对所有 $t\in[0,T)$。于是，当 $t$ 充分小时，
\begin{align}
\Omega_{t}=\{y\in\mathbb{R}^{3}:y=x+tv_{0}(x),\ x\in\Omega\},
\label{align:ch04-Omega_t}
\end{align}
由此
\begin{align}
|\Omega_{t}|=\int_{\Omega}\det[I+t\operatorname{grad}v_{0}(x)]\,\mathrm{d}x.
\label{align:ch04-体积积分}
\end{align}
因而，当 $t$ 小时，$|\Omega_{t}|$ 是多项式：
\begin{align}
|\Omega_{t}|=Jt^{3}+Kt^{2}+Lt+N,
\label{align:ch04-体积多项式}
\end{align}
其中 $J,K,L$ 依赖于 $\operatorname{grad}v_{0}$，而 $N=|\Omega|$。实际上，常数 $J,K,L$ 是 $v_{0}$ 的零 Lagrangian 在 $\Omega$ 上的积分，因而完全由 $v_{0}$ 与 $\operatorname{grad}v_{0}$ 在 $\Gamma$ 上的值确定。特别地，
\begin{align}
J=\int_{\Omega}\det[\operatorname{grad}v_{0}(x)]\,\mathrm{d}x.
\label{align:ch04-J}
\end{align}

由于上述论证在把 $0$ 替换为任一 $\bar{t}\in[0,T)$ 后可重复，故\eqref{align:ch04-体积多项式}对所有 $t\in[0,T)$ 成立。现在显然，若 $J\leqslant 0$，则解无法在 $t$ 大时满足 $|\Omega_{t}|\geqslant At^{3}$。因而，当初值满足 $J\leqslant 0$ 时，任一古典解的寿命必有限。

在当前情形下，古典解的消亡既可作为波破碎的结果，也可作为气--真空界面间形成奇点的结果。让我们在下列设定中讨论后一可能性：界面 $\Gamma_{t}$ 仍为光滑曲面，解 $(\rho,v)$ 在 $\Omega_{t}$ 上仍光滑，但允许在 $\Gamma_{t}$ 上奇异。

首先检验界面是否可为激波。为此，我们以 $v$ 为 $\Omega_{t}$ 的单位外法向应用 Rankine--Hugoniot 条件\cref{chapter:ch03-双曲平衡律方程组}（3.3.40），推出压力 $p$ 的跳跃、从而密度 $\rho$ 的跳跃为零。因而界面处的奇点必弱于跳跃间断。我们寻求允许界面以受控速率加速的条件。自然的假设是声速 $c$ 的平方 Lipschitz，且其法向导数跨越界面有跳跃：

\begin{align}
\left[\left[\frac{\partial c^{2}}{\partial v}\right]\right]=g,
\label{align:ch04-界面加速度}
\end{align}
其中 $g$ 为有界函数，且必为非负。事实上，当\eqref{align:ch04-界面加速度}成立时，方程组\eqref{align:ch04-Euler对称双曲}第二式限制到界面上给出
\begin{align}
\partial_{t}v+(v\cdot\operatorname{grad})v=\frac{g}{\gamma-1}\nu,
\label{align:ch04-界面加速度2}
\end{align}
应与\eqref{align:ch04-界面输运}对比。由于 $g\geqslant 0$，\eqref{align:ch04-界面加速度2}标志着界面的加速。特别地，当 $g$ 充分大时，$\Omega_{t}$ 的体积可以与要求 $|\Omega_{t}|\geqslant At^{3}$ 相容的速率增长，从而允许在界面上真空奇点之外无其他奇点的整体解存在。更多见\cref{chapter:ch05-熵与经典解稳定性}。


我们现转向古典解因波破碎而崩溃的问题。正如在第 4.2 小节 Burgers 方程情形所见，这一现象的物理描述是：由于非线性，波速依赖于波幅，这可能产生陡峭波剖面并最终破碎。同一机制在 Euler 方程组中也存在。事实上，正如我们将在\cref{chapter:ch07-一维双曲平衡律组}中看到，一维空间中 Euler 方程组的压缩波必破碎。然而在三维中情形更微妙。一开始，压缩所致波剖面变陡与色散相竞争，后者起相反作用。这延迟、甚至可能挫败波的破碎。困难还因涡量的存在而加重，其对奇点发展的潜在贡献尚未被完全理解。因而为专注于压缩效应，我们通过把讨论限于无旋流动来消除涡量的影响。

在无旋流动中，速度由位势 $\phi$ 经\cref{chapter:ch03-双曲平衡律方程组}（3.3.48）导出，Euler 方程组化为\cref{chapter:ch03-双曲平衡律方程组}（3.3.49）、（3.3.50），其中 $h$ 为焓。为简单起见设体积力为零，$g=0$。由于 $h'(\rho)=p'(\rho)/\rho>0$，可把密度实现为焓的函数 $\rho=\rho(h)$。


下文中我们采用求和约定。把作为 $h$ 函数的 $\rho$（其中 $h$ 由\cref{chapter:ch03-双曲平衡律方程组}（3.3.49）给出）代入质量守恒方程，化为准线性波动方程

\begin{align}
\partial_{t}^{2}\phi-c^{2}(\rho)\Delta\phi=-2(\partial_{j}\phi)(\partial_{t}\partial_{j}\phi)-(\partial_{i}\phi)(\partial_{j}\phi)(\partial_{i}\partial_{j}\phi),
\label{align:ch04-准线性波动}
\end{align}
其中 $\phi$ 为位势，$c(\rho)$ 表示声速。注意，在任一静止状态 $\rho=\rho_{0}$、$v=0$ 处对\eqref{align:ch04-准线性波动}线性化，给出古典波动方程。

\eqref{align:ch04-准线性波动}一个显著特征是它可实现为 Euler--Lagrange 方程。事实上，回顾 $\partial_{i}\phi=v_{i}$（$i=1,2,3$），并令 $\partial_{t}\phi=-v_{0}$。引入 Lagrangian
\begin{align}
L(v_{0},v_{1},v_{2},v_{3})=p(\rho(h)),\qquad h=v_{0}-\tfrac{1}{2}v_{i}v_{i}
\label{align:ch04-Lagrangian}
\end{align}
并利用 $\frac{\mathrm{d}p}{\mathrm{d}\rho}\frac{\mathrm{d}\rho}{\mathrm{d}h}=\rho$，推出
\begin{align}
\frac{\partial L}{\partial v_{0}}=\rho,\qquad \frac{\partial L}{\partial v_{i}}=-\rho v_{i},\qquad i=1,2,3,
\label{align:ch04-L偏导}
\end{align}
这蕴含
\begin{align}
\partial_{t}\frac{\partial L}{\partial v_{0}}-\partial_{i}\frac{\partial L}{\partial v_{i}}=\partial_{t}\rho+\partial_{i}(\rho v_{i})=0.
\label{align:ch04-EulerLagrange}
\end{align}

我们考虑\eqref{align:ch04-准线性波动}在指定初值
\begin{align}
\phi(x,0)=\phi_{0}(x),\qquad \partial_{t}\phi(x,0)=\phi_{1}(x),\qquad x\in\mathbb{R}^{3}
\label{align:ch04-波动初值}
\end{align}
下的 Cauchy 问题。作为\cref{chapter:ch05-熵与经典解稳定性}（第 V 章将证明）的推论，每当 $\phi_{0}\in W^{4,2}(\mathbb{R}^{3})$、$\phi_{1}\in W^{3,2}(\mathbb{R}^{3})$，存在唯一的古典 $C^{2}$ 解 $\phi$，局部在时间内。人们已付出大量努力在初值属于足够高阶 Sobolev 空间、范数大小由参数 $\varepsilon$ 界定的假设下监测古典解。这一任务需要强有力的分析，充分利用\eqref{align:ch04-准线性波动}的内在几何结构，特别是其作为 Euler--Lagrange 方程的体现。不幸的是，这项工作过于技术化、过于费力，即使以节略形式也无法在此呈现，读者应查阅第 4.9 小节所引文献。下一段给出主要结论的梗概。


在三维中，色散对古典波动方程 Cauchy 解的导数诱导 $O(t^{-1})$ 衰减率。\eqref{align:ch04-准线性波动}的解在一阶导数层面继承这一性质。然而在二阶导数层面，非线性与色散之间势均力敌。色散成功地把古典解的寿命延长到 $O(\exp\frac{1}{\varepsilon})$，相对于一维情形（寿命仅为 $O(\frac{1}{\varepsilon})$）是重大改进。尽管如此，非线性最终占优，驱使二阶导数趋于无穷。富有洞察力的证明通过引入适配波剖面的特殊坐标来进行：识别二阶导数沿之增长并最终破碎的主方向，与之相对的是色散占优、控制导数大小的横截方向。这些坐标下的分析相当显式，因而除了暴露波的破碎，它还给出解最大发展的描述。

我们以关于 Euler 方程组弱解的某些惊人、也许令人不安的事实结束本节。即将到来的第 5.2 小节的分析将确立：每当\eqref{align:ch04-Euler}的 Cauchy 问题在某时间区间 $[0,T)$ 上有古典解，此解是唯一的--不仅在其他古典解之中，甚至在相对于熵--熵流对\eqref{align:ch04-Euler熵}满足熵可容性条件\eqref{align:ch04-熵积分0}的更广 $L^{\infty}$ 弱解类之中。然而，在缺乏古典解时，上述熵不等式不再足以挑选出唯一可容解。这一显著事实将通过考虑下列 Cauchy 问题来展示：初值在 $x$ 的每个分量 $x_{i}$ 上以周期 1 为周期。因而把时刻 $t$ 的解视为定义在标准环面 $\mathbb{T}^{3}$ 上的函数将是方便的。多解的存在性由下述定理确立。

\begin{theorem}\label{theorem:ch04-Euler无穷弱解}
存在 $m_{0}\in L^{\infty}(\mathbb{T}^{3})$，使得 Euler 方程组\eqref{align:ch04-Euler}以初值
\begin{align}
\rho(x,0)=1,\qquad m(x,0)=m_{0}(x),\qquad x\in\mathbb{T}^{3}
\label{align:ch04-Euler环面初值}
\end{align}
在 $[0,\infty)$ 上有无穷多个 $L^{\infty}$ 弱解 $(\rho,m)$，它们相对于熵--熵流对\eqref{align:ch04-Euler熵}满足可容性条件\eqref{align:ch04-熵积分tau}。
\end{theorem}

证明冗长且技术化，这里仅作粗略勾勒。细节见第 4.9 小节所引文献。

Euler 方程组在 $L^{\infty}$ 弱解设定下一个惊人特征是：Cauchy 问题是高度欠定的，以至于可同时指定密度场 $\rho$ 与动量场长度 $|m|$，却仍留下构造无穷多个满足熵可容性条件的解的余地。于是，令 $\rho\equiv 1$，此时\eqref{align:ch04-Euler}化为
\begin{align}
\begin{cases}
\operatorname{div}m^{\top}=0\\
\partial_{t}m+\operatorname{div}(mm^{\top})=0.
\end{cases}
\label{align:ch04-密度归一}
\end{align}

暂时搁置熵可容性要求，固定 $T>0$ 与一个在 $[0,T]$ 上取正值的 $C^{1}$ 函数 $M$，在 $[0,T]$ 上求\eqref{align:ch04-密度归一}的 $L^{\infty}$ 解，使其满足初值与终值条件
\begin{align}
m(\cdot,0)=m_{0}(\cdot),\qquad m(\cdot,T)=0,
\label{align:ch04-两端条件}
\end{align}
以及约束
\begin{align}
|m(x,t)|=M(t),\qquad \text{a.e. 在 }\mathbb{T}^{3}\times(0,T)\text{ 上}.
\label{align:ch04-模约束}
\end{align}
显然的相容条件是 $|m_{0}(x)|\leqslant M(0)$，a.e. 在 $\mathbb{T}^{3}$ 上。

设\eqref{align:ch04-模约束}成立，可把\eqref{align:ch04-密度归一}重写为等价形式
\begin{align}
\begin{cases}
\operatorname{div}m^{\top}(x,t)=0\\
\partial_{t}m(x,t)+\operatorname{div}U(x,t)=0,
\end{cases}
\label{align:ch04-密度归一2}
\end{align}
\begin{align}
U(x,t)=m(x,t)m^{\top}(x,t)-\tfrac{1}{3}|m(x,t)|^{2}I.
\label{align:ch04-U定义}
\end{align}

于是，按连续物理的精神，我们把\eqref{align:ch04-密度归一}视为守恒律组\eqref{align:ch04-密度归一2}与本构方程\eqref{align:ch04-U定义}的复合。注意 $3\times 3$ 矩阵值函数 $U$ 对称且无迹。下文 $\mathbb{N}$ 表示对称无迹 $3\times 3$ 矩阵空间。

分析的第一步是引入一类称为上述方程组\emph{次解}的函数。

\eqref{align:ch04-两端条件}、\eqref{align:ch04-模约束}、\eqref{align:ch04-密度归一2}与\eqref{align:ch04-U定义}的次解是一个定义在 $\mathbb{T}^{3}\times[0,T]$ 上、取值于 $\mathbb{R}^{3}$ 的函数 $u$，具有下列性质：
\begin{enumerate}
\item\label{item:ch04-次解-a} 函数 $t\mapsto u(\cdot,t)$ 在 $[0,T]$ 上于 $L^{\infty}(\mathbb{T}^{3};\mathbb{R}^{3})$ weak$^{*}$ 下连续，且
\begin{align}
u(\cdot,0)=m_{0}(\cdot),\qquad u(\cdot,T)=0.
\label{align:ch04-次解边界}
\end{align}
\item\label{item:ch04-次解-b} $u$ 在 $\mathbb{T}^{3}\times(0,T)$ 上为 $C^{1}$，且满足
\begin{align}
\begin{cases}
\operatorname{div}u^{\top}(x,t)=0\\
\partial_{t}u(x,t)+\operatorname{div}U(x,t)=0,
\end{cases}
\label{align:ch04-次解方程}
\end{align}
其中 $U$ 是某在 $\mathbb{T}^{3}\times(0,T)$ 上取值于 $\mathbb{N}$ 的 $C^{1}$ 函数。
\item\label{item:ch04-次解-c} 对任一 $x\in\mathbb{T}^{3}$ 与 $t\in(0,T)$，
\begin{align}
u(x,t)u^{\top}(x,t)-U(x,t)\leqslant\tfrac{1}{3}M^{2}(t)I.
\label{align:ch04-次解凸条件}
\end{align}
\end{enumerate}

应注意，对任一固定 $w\in\mathbb{R}^{3}$ 与 $W\in\mathbb{N}$，
\begin{align}
ww^{\top}-W\geqslant\tfrac{1}{3}|w|^{2}I.
\label{align:ch04-凸性恒等式}
\end{align}
因而\eqref{align:ch04-次解凸条件}蕴含 $|u(x,t)|\leqslant M(t)$，对所有 $x\in\mathbb{T}^{3}$ 与 $t\in(0,T)$。

我们在与 $m_{0}$、$M$、$T$ 相联系的次解集非空的假设下进行。记此集合为 $\mathscr{X}$。特别地，当 $m_{0}\equiv 0$ 时，对任一正的 $M$，与之相联系的 $\mathscr{X}$ 非空，因为它含有零函数。记 $\overline{\mathscr{X}}$ 为 $\mathscr{X}$ 在 $L^{2}(\mathbb{T}^{3}\times[0,T])$ 弱拓扑下的闭包。

在 $\overline{\mathscr{X}}$ 上定义泛函
\begin{align}
J[u]=\int_{0}^{T}\int_{\mathbb{T}^{3}}\left[|u(x,t)|^{2}-M^{2}(t)\right]\,\mathrm{d}x\mathrm{d}t,
\label{align:ch04-J泛函}
\end{align}
其值非正。$\overline{\mathscr{X}}$ 的成员在分布意义下继承 $\mathscr{X}$ 的性质\eqref{align:ch04-次解方程}与\eqref{align:ch04-次解凸条件}，几乎处处成立，因为对称矩阵的最大特征值是凸函数。由这一观察与\eqref{align:ch04-凸性恒等式}可知：$m\in\overline{\mathscr{X}}$ 是\eqref{align:ch04-两端条件}、\eqref{align:ch04-模约束}、\eqref{align:ch04-密度归一2}与\eqref{align:ch04-U定义}的解当且仅当 $J[m]=0$。

由于对所有 $u\in\mathscr{X}$ 有 $|u(x,t)|\leqslant M(t)$，故任一满足 $J[m]=0$ 的 $m\in\overline{\mathscr{X}}$ 必为 $J$ 的连续点。事实证明逆命题也成立：若 $J$ 在 $m\in\overline{\mathscr{X}}$ 处连续，则 $J[m]=0$。理由是：若对某 $u\in\mathscr{X}$ 有 $J[u]<0$，则存在 $v\in\mathscr{X}$ 与 $u$ 弱接近，但 $J[v]$ 与 $J[u]$ 相差甚大。精确命题由下述引理给出，其技术证明见第 4.9 小节所引文献。

\begin{lemma}\label{lemma:ch04-次解扰动}
设 $u$ 为任一次解，$(\tau_{1},\tau_{2})$ 为 $(0,T)$ 的任一子区间。则存在序列 $\{v_{n}\}\subset\mathscr{X}$，使得 $u-v_{n}$ 支于时间区间 $(\tau_{1},\tau_{2})$ 内，$v_{n}\to u$ 在 $L^{\infty}$ weak$^{*}$ 下，且
\begin{align}
\liminf_{n\to\infty}\int_{0}^{T}\int_{\mathbb{T}^{3}}|v_{n}(x,t)-u(x,t)|^{2}\,\mathrm{d}x\mathrm{d}t\geqslant a\int_{\tau_{1}}^{\tau_{2}}\int_{\mathbb{T}^{3}}\left[|u(x,t)|^{2}-M^{2}(t)\right]^{2}\,\mathrm{d}x\mathrm{d}t
\label{align:ch04-次解扰动估计}
\end{align}
成立，其中 $a>0$ 与 $u$ 无关。
\end{lemma}

次解 $u$ 在 $\mathscr{X}$ 内的扰动--在弱拓扑意义下“微小”、却对熵引起显著变化--典型地涉及快速振荡。事实上，剧烈振荡的存在正是本理论的标志。

我们把 $\overline{\mathscr{X}}$ 视为完备度量空间。泛函 $J$ 属 Baire 1 类（连续函数的逐点极限），因而其连续点集在 $\overline{\mathscr{X}}$ 中稠密。这些点每一个都给出\eqref{align:ch04-密度归一}、\eqref{align:ch04-两端条件}、\eqref{align:ch04-模约束}的解。此后我们称这些解为\emph{奇异}（exotic）解。

我们现转向奇异解的可容性问题。由\eqref{align:ch04-Euler非守恒}与\eqref{align:ch04-模约束}，
\begin{align}
\eta(x,t)=\tfrac{1}{2}M^{2}(t)+\varepsilon(1),\qquad Q(x,t)=\left[\tfrac{1}{2}M^{2}(t)+\varepsilon(1)+p(1)\right]m(x,t).
\label{align:ch04-奇异熵}
\end{align}
注意 $Q^{\top}$ 的散度为零。因而当 $\dot{M}(t)\leqslant 0$ 时，解 $(1,m)$ 在 $\mathbb{T}^{3}\times(0,\infty)$ 上满足熵可容性判据\eqref{align:ch04-熵不等式}。特别地，若 $M$ 为常数，则熵在 $\mathbb{T}^{3}\times(0,T)$ 上守恒。

上述性质突显了奇异解与标准解（例如 $BV$ 类）之间的差异。在后一情形，熵仅由跳跃间断产生。正如在（子）节 3.3.6 中注意到，Euler 方程组的解可支承两类跳跃间断，即压缩激波与涡面（接触间断）。产生熵的压缩激波不能参与奇异解，因为它们涉及密度跳跃。相反，与均匀密度相容的涡面很可能出现在奇异解中，并可作为这些解振荡剖面的构造单元。然而，涡面不产生熵，因而\eqref{align:ch04-奇异熵}中编码的 $\eta$ 的衰减不能归因于跳跃间断的存在，而来自另一机制。我们将在\cref{chapter:ch06-标量守恒律L1理论}与\cref{chapter:ch11-一维标量守恒律}中讨论相关问题。


尽管 $\dot{M}(t)\leqslant 0$ 保证奇异解满足\eqref{align:ch04-熵不等式}，更具选择性的熵可容性判据\eqref{align:ch04-熵积分tau}并不成立，除非 $t=0$ 是函数 $t\mapsto m(\cdot,t)$ 在强 $L^{1}(\mathbb{T}^{3};\mathbb{R}^{3})$ 拓扑下的右连续点。事实证明，给定 $m\in\mathscr{X}$，可借助\cref{lemma:ch04-次解扰动}中编码的 $\mathscr{X}$ 性质构造 $\bar{m}\in\mathscr{X}$，使得对某个 $\tau\in(0,T)$ 有 $\bar{m}(\cdot,\tau)=M(\tau)$，而对 $t$ 接近 $T$ 有 $\bar{m}(\cdot,t)=m(\cdot,t)$。现在显然，若以 $\bar{M}$ 替换 $M$、以 $\bar{m}_{0}$ 替换 $m_{0}$，其中 $\bar{M}(t)=M(t+\tau)$，$\bar{m}_{0}(\cdot)=\bar{m}(\cdot,\tau)$，所得奇异解将满足\eqref{align:ch04-熵积分tau}。

精细的技术分析表明，事实上甚至可构造耗散熵的奇异解，它们是 Hölder（而非 Lipschitz）连续的。

Euler 方程组具有大规模非唯一性的奇异解这一发现令人震惊，因而人们希望它们会因物理上不可容而被排除。尽管如此，我们上面看到这些解通过了熵可容性检验。另一方面，由于在设定 $M(t)$ 时所获得的灵活性，没有一个奇异解会满足极大耗散性判据\eqref{align:ch04-极大耗散}。另一风险是：先前被判定为极大耗散的标准解，在与奇异解的竞争中可能失去这一地位。读者可在\cref{chapter:ch09-容许波扇与riemann问题}中找到关于这一问题的相关评论。


\end{document}
